%% ==========================================================================
%%  4  Obstructions to the Natural Approaches
%%
%%  MAIN REPORT ONLY -- not included by working.tex, which carries the full
%%  approach sections.  Every label here is prefixed `main:' so that nothing
%%  can collide with working/.
%%
%%  Covers the first two obstructions only.  The third needs the replica
%%  instance and is therefore deferred to sections/replica.tex.
%%
%%  Condensed from working/approach_1.tex (thm:freeinsert, thm:cyclebound,
%%  thm:obstruction) and working/approach_2.tex (prop:msw-false).
%% ==========================================================================
\section{Obstructions to the Natural Approaches}
\label{sec:obstructions}

Three techniques settle the neighbouring corners of the square in
\Cref{sec:intro}: build the allocation one item at a time, as in~\cite{BKNS22};
select a welfare-maximising allocation, as the exchange structure of dichotomous
valuations invites; or reduce to a goods instance and price the residual
constraint away. Each fails for negative dichotomous valuations, and each failure
is witnessed by a single small instance. We treat the first two here; the third
needs a construction of its own and is deferred to \Cref{sec:replica}. These
obstructions delimit what a proof of Conjecture~\ref{conj:main} can look like,
and two of the tools built along the way are used throughout.

% --------------------------------------------------------------------------
\subsection{Incremental insertion}
\label{sec:main-insertion}

The algorithm of~\cite{BKNS22} maintains an envy-free solution over a partial
allocation with $\subsidy \in \set{0,1}^n$, allocates one further item at each
step, and repairs the invariant by permuting whole bundles. Nothing already
allocated is ever moved. Two steps of that template come out \emph{easier} on
the chore side.

\begin{theorem}[Free insertion]
\label{main:freeinsert}
Let $\alloc$ be envy-freeable and let $g$ be an unallocated chore with
$\cost_x(\bundle{x} \cup \set{g}) = \cost_x(\bundle{x})$ for some agent $x$.
Then the allocation obtained by giving $g$ to $x$ is envy-freeable, and its
subsidy vector is pointwise no larger than that of $\alloc$.
\end{theorem}

\begin{proof}
By Lemma~\ref{lem:arcupdate} the arcs between agents other than $x$ are
unchanged; the arcs out of $x$ are unchanged, the marginal being $0$; and the
arcs into $x$ weakly decrease. Every arc weight therefore weakly decreases, so
no positive-weight cycle can appear and no path weight can rise.
\end{proof}

For goods the corresponding step needs the receiver to be maximally subsidised,
a permutation of the bundles, and an extendability argument, because there the
insertion \emph{raises} the arcs into the receiver. Here it lowers them. The
only case requiring work is a chore whose marginal cost is $1$ for every agent
on her own bundle.

The second tool is the cheapest general handle on $\pathw{\alloc}$ available,
and is used repeatedly below.

\begin{theorem}[Cycle-closing bound]
\label{main:cyclebound}
For any envy-freeable allocation $\alloc$ and any agent $u$,
\[
  \pathw{\alloc}(u) \;\le\; \max\Big(0, \ \max_{v \ne u}
  \big[\,\cost_v(\bundle{u}) - \cost_v(\bundle{v})\,\big]\Big).
\]
\end{theorem}

\begin{proof}
Since $\alloc$ is envy-freeable, $\envygraph{\alloc}$ has no positive-weight
directed cycle (Theorem~\ref{thm:hs-characterisation}(iii)). Every closed walk
decomposes into simple cycles and so also has weight at most $0$; deleting a
closed sub-walk therefore never decreases a walk's weight, and every walk
starting at $u$ is dominated by a \emph{simple} path starting at $u$. As there
are finitely many simple paths, the maximum defining $\pathw{\alloc}(u)$ is
attained, by a simple path $P^{\ast}$.

If $P^{\ast}$ is empty then $\pathw{\alloc}(u) = 0$ and the claim holds, the
right-hand side being non-negative. Otherwise $P^{\ast}$, being simple, does not
revisit $u$, so its terminal vertex $v$ satisfies $v \ne u$ and appending the arc
$(v,u)$ yields a simple cycle $C$ with $\edgew{\alloc}(C) \le 0$. Hence
\[
  \pathw{\alloc}(u) \;=\; \edgew{\alloc}(P^{\ast})
  \;=\; \edgew{\alloc}(C) - \edgew{\alloc}(v,u)
  \;\le\; -\edgew{\alloc}(v,u)
  \;=\; \cost_v(\bundle{u}) - \cost_v(\bundle{v}),
\]
which is at most the maximum over $v \ne u$.
\end{proof}

\begin{corollary}
\label{main:onestep}
If $\cost_v(\bundle{u}) \le \cost_v(\bundle{v}) + 1$ for all $u,v \in \agents$
--- no agent would suffer more than one extra unit by taking anybody else's
bundle --- then $\optsubsidy \in \set{0,1}^n$.
\end{corollary}

The bound is an \emph{upper estimate}, not an evaluation. Here
$\pathw{\alloc}(u)$ is the maximum weight of a path out of $u$, whereas the
right-hand side retains only the endpoint of that path, and the inequality is in
general strict: if $u \to a \to b$ has weight $2$ and
$\edgew{\alloc}(b,u) = -3$, the cycle weighs $-1$, so $\alloc$ is envy-freeable
with $\pathw{\alloc}(u) = 2$ while the bound returns $3$. Evaluating
$\pathw{\alloc}$ requires a longest-path computation; the bound costs $O(n)$ arc
lookups and is used only where an upper estimate suffices, as in
Corollary~\ref{main:onestep}.

The template nevertheless cannot be completed.

\begin{example}[Insertion cannot be completed]
\label{main:ex-insertion}
Take $\items = \set{a_1,a_2,g}$ and
\[
  \cost_1(S) = \max\big(0, \abs{S} - 1\big), \qquad
  \cost_2(S) = \cost_3(S) = \abs{S},
\]
all dichotomous. The partial allocation
$\alloc = (\set{a_1,a_2}, \emptyset, \emptyset)$ has envy matrix
\[
  \big(\edgew{\alloc}(i,j)\big) =
  \begin{pmatrix} 0 & 1 & 1 \\ -2 & 0 & 0 \\ -2 & 0 & 0 \end{pmatrix},
\]
with no positive cycle and $\optsubsidy = (1,0,0)$: a legal state of the
invariant. The marginal cost of $g$ is $1$ for every agent on her own bundle, and
inserting $g$ into each bundle in turn gives minimum subsidy vectors
$(2,0,0)$, $(2,1,0)$ and $(2,0,1)$ --- each with an entry equal to $2$. Yet the
complete allocation $(\set{a_1}, \set{a_2}, \set{g})$ has an identically zero
envy matrix and needs no subsidy at all.
\end{example}

Note that $\cost_1$ is supermodular, which is permitted: neither~\cite{BKNS22}
nor the present work restricts to submodular valuations, and the instance
genuinely needs it. The conclusion is that any proof of
Conjecture~\ref{conj:main} must be free to \textbf{re-open already-allocated
bundles}; repairing by permutation alone is provably insufficient.

% --------------------------------------------------------------------------
\subsection{Selecting a utilitarian optimum}
\label{sec:main-utilitarian}

If the allocation cannot be built incrementally, it may instead be
\emph{selected} from a set small enough to reason about. Utilitarian optimality
is the obvious candidate, for two good reasons. By
Theorem~\ref{thm:hs-characterisation} an allocation minimising
$\sum_i \cost_i(\bundle{i})$ over \emph{all} allocations is automatically
envy-freeable, so half of what is needed comes free. And it carries an exchange
property: if $g \in \bundle{i}$ satisfies $\cost_i(\bundle{i}) -
\cost_i(\bundle{i} \setminus \set{g}) = 1$, then $\cost_j(\bundle{j} \cup
\set{g}) = \cost_j(\bundle{j}) + 1$ for every $j$, since otherwise the
allocation would not have been optimal.

\begin{example}[The utilitarian optimum can be strictly worse]
\label{main:ex-utilitarian}
Take $\items = \set{g_1,g_2,g_3,g_4}$ with the costs below; each column is $0$ at
$\emptyset$ with every marginal in $\set{0,1}$, so all three are dichotomous.

\begin{center}\small
\begin{tabular}{@{}lccc@{\hspace{2.5em}}lccc@{}}
\toprule
$S$ & $\cost_1$ & $\cost_2$ & $\cost_3$ &
$S$ & $\cost_1$ & $\cost_2$ & $\cost_3$ \\
\midrule
$\emptyset$      & 0 & 0 & 0 & $g_3g_4$       & 2 & 1 & 1 \\
$g_1$            & 1 & 1 & 1 & $g_1g_2g_3$    & 2 & 3 & 2 \\
$g_2$            & 1 & 1 & 1 & $g_1g_2g_4$    & 3 & 2 & 2 \\
$g_3$            & 1 & 1 & 0 & $g_1g_3g_4$    & 2 & 2 & 2 \\
$g_4$            & 1 & 1 & 1 & $g_2g_3g_4$    & 2 & 2 & 2 \\
$g_1g_2$         & 2 & 2 & 2 & $g_1g_2g_3g_4$ & 3 & 3 & 3 \\
$g_1g_3$         & 1 & 2 & 1 & & & & \\
$g_1g_4$         & 2 & 2 & 2 & & & & \\
$g_2g_3$         & 1 & 2 & 1 & & & & \\
$g_2g_4$         & 2 & 2 & 2 & & & & \\
\bottomrule
\end{tabular}
\end{center}

Agent~1 pays $1$ for every singleton and agent~3 pays $0$ only for
$\set{g_3}$, so any allocation of total cost below $2$ would force
$\bundle{2} \supseteq \set{g_1,g_2,g_4}$, costing agent~2 at least $2$. The
minimum total cost is therefore $2$, attained \emph{only} by
$\alloc^{\star} = (\emptyset,\ \set{g_1,g_2,g_4},\ \set{g_3})$, whose envy graph
has $\edgew{\alloc^\star}(2,1) = 2$ and hence $\optsubsidy = (0,2,0)$. But
$\allocB = (\set{g_1,g_3},\ \set{g_2},\ \set{g_4})$, of total cost $3$, has every
arc at most $0$ and needs no subsidy.
\end{example}

The optimum here is \emph{unique}, so there is no tie to break: every rule
selecting a nonempty subset of the utilitarian-optimal allocations returns
$\alloc^{\star}$ and pays $2$, while a subsidy-free allocation exists outside
that set. The gap is the largest possible in one step, for one extra unit of
total burden. Any search procedure must therefore be free to \textbf{give up
total cost}.
